
\documentclass{article}
\usepackage{colortbl}
\usepackage{makecell}
\usepackage{multirow}
\usepackage{supertabular}

\begin{document}

\newcounter{utterance}

\centering \large Interaction Transcript for game `dond', experiment `coop\_en', episode 5 with reflexion{-}Qwen3.5{-}4B.
\vspace{24pt}

{ \footnotesize  \setcounter{utterance}{1}
\setlength{\tabcolsep}{0pt}
\begin{supertabular}{c@{$\;$}|p{.15\linewidth}@{}p{.15\linewidth}p{.15\linewidth}p{.15\linewidth}p{.15\linewidth}p{.15\linewidth}}
    \# & $\;$A & \multicolumn{4}{c}{Game Master} & $\;\:$B\\
    \hline

    \theutterance \stepcounter{utterance}  
    & & \multicolumn{4}{p{0.6\linewidth}}{
        \cellcolor[rgb]{0.9,0.9,0.9}{
            \makecell[{{p{\linewidth}}}]{
                \texttt{\tiny{[P1$\langle$GM]}}
                \texttt{You are playing a cooperative negotiation game in which you have to agree on how to divide a set of items among you and another player.} \\
\\ 
\texttt{Rules:} \\
\texttt{(a) You and the other player are given a set of items. Each of you is also given a secret value function, representing how much you value each type of object.} \\
\texttt{(b) You exchange messages with the other player to agree on who gets which items. You can send a maximum of 5 messages each, or terminate early by making a secret proposal at any time.} \\
\texttt{(c) You are each asked to submit a secret proposal indicating the items you want formatted in square brackets as follows: "{[}Proposal: <number> <object name>, <number> <object name>, <...>{]}"} \\
\texttt{(d) If your proposals are complementary, i.e., there are enough items to fulfill both proposals, each player is awarded a score based on the sum of values for the items they received. Otherwise, both of you get zero points.} \\
\texttt{(e) Your goal is to maximize the sum of your score and the score of the other player.} \\
\\ 
\texttt{Let us start.} \\
\\ 
\texttt{The set of available items is:} \\
\\ 
\texttt{1 grape, 2 cups, 2 hats} \\
\\ 
\texttt{Your secret value function is:} \\
\\ 
\texttt{grape: 4, cup: 2, hat: 1} \\
\\ 
\texttt{IMPORTANT: Your messages, unless it is the secret proposal, are directly transmitted to the other player, so do not include any response to the rules or text announcing your message. To make a secret proposal, use the indicated format. Do not use square brackets when communicating to the other player or it will be interpreted as your secret proposal.} \\
            }
        }
    }
    & \\ \\

    \theutterance \stepcounter{utterance}  
    & \multicolumn{4}{p{0.6\linewidth}}{
        \cellcolor[rgb]{0.9,0.9,0.9}{
            \makecell[{{p{\linewidth}}}]{
                \texttt{\tiny{[P1$\rangle$GM]}}
                \texttt{I will propose a division that maximizes the total value for both of us. Since I value grapes the most (4), I should keep the grape. The cups are worth 2 each, and hats are worth 1 each.} \\
\\ 
\texttt{To maximize the sum of our scores, I will propose that I take the grape and one cup (total value 6), and you take the remaining cup and both hats (total value 4). This ensures both of us get items, avoiding the zero{-}point penalty, and distributes the high{-}value items efficiently.} \\
\\ 
\texttt{Cup: 2, hat: 1} \\
            }
        }
    }
    & & \\ \\

    \theutterance \stepcounter{utterance}  
    & & & \multicolumn{2}{p{0.3\linewidth}}{
        \cellcolor[rgb]{0.9,0.9,0.9}{
            \makecell[{{p{\linewidth}}}]{
                \texttt{\tiny{[GM$|$GM]}}
                \texttt{continue} \\
            }
        }
    }
    & & \\ \\

    \theutterance \stepcounter{utterance}  
    & & & \multicolumn{2}{p{0.3\linewidth}}{
        \cellcolor[rgb]{0.9,0.9,0.9}{
            \makecell[{{p{\linewidth}}}]{
                \texttt{\tiny{[GM$|$GM]}}
                \texttt{I will propose a division that maximizes the total value for both of us. Since I value grapes the most (4), I should keep the grape. The cups are worth 2 each, and hats are worth 1 each.} \\
\\ 
\texttt{To maximize the sum of our scores, I will propose that I take the grape and one cup (total value 6), and you take the remaining cup and both hats (total value 4). This ensures both of us get items, avoiding the zero{-}point penalty, and distributes the high{-}value items efficiently.} \\
\\ 
\texttt{Cup: 2, hat: 1} \\
            }
        }
    }
    & & \\ \\

    \theutterance \stepcounter{utterance}  
    & & \multicolumn{4}{p{0.6\linewidth}}{
        \cellcolor[rgb]{0.9,0.9,0.9}{
            \makecell[{{p{\linewidth}}}]{
                \texttt{\tiny{[GM$\rangle$P2]}}
                \texttt{You are playing a cooperative negotiation game in which you have to agree on how to divide a set of items among you and another player.} \\
\\ 
\texttt{Rules:} \\
\texttt{(a) You and the other player are given a set of items. Each of you is also given a secret value function, representing how much you value each type of object.} \\
\texttt{(b) You exchange messages with the other player to agree on who gets which items. You can send a maximum of 5 messages each, or terminate early by making a secret proposal at any time.} \\
\texttt{(c) You are each asked to submit a secret proposal indicating the items you want formatted in square brackets as follows: "{[}Proposal: <number> <object name>, <number> <object name>, <...>{]}"} \\
\texttt{(d) If your proposals are complementary, i.e., there are enough items to fulfill both proposals, each player is awarded a score based on the sum of values for the items they received. Otherwise, both of you get zero points.} \\
\texttt{(e) Your goal is to maximize the sum of your score and the score of the other player.} \\
\\ 
\texttt{Let us start.} \\
\\ 
\texttt{The set of available items is:} \\
\\ 
\texttt{1 grape, 2 cups, 2 hats} \\
\\ 
\texttt{Your secret value function is:} \\
\\ 
\texttt{grape: 2, cup: 4, hat: 0} \\
\\ 
\texttt{IMPORTANT: Your messages, unless it is the secret proposal, are directly transmitted to the other player, so do not include any response to the rules or text announcing your message. To make a secret proposal, use the indicated format. Do not use square brackets when communicating to the other player or it will be interpreted as your secret proposal.} \\
\\ 
\\ 
\texttt{I will propose a division that maximizes the total value for both of us. Since I value grapes the most (4), I should keep the grape. The cups are worth 2 each, and hats are worth 1 each.} \\
\\ 
\texttt{To maximize the sum of our scores, I will propose that I take the grape and one cup (total value 6), and you take the remaining cup and both hats (total value 4). This ensures both of us get items, avoiding the zero{-}point penalty, and distributes the high{-}value items efficiently.} \\
\\ 
\texttt{Cup: 2, hat: 1} \\
            }
        }
    }
    & \\ \\

    \theutterance \stepcounter{utterance}  
    & & & \multicolumn{4}{p{0.6\linewidth}}{
        \cellcolor[rgb]{0.9,0.9,0.9}{
            \makecell[{{p{\linewidth}}}]{
                \texttt{\tiny{[GM$\langle$P2]}}
                \texttt{I will propose a division that maximizes the total value for both of us. Since I value cups the most (4), I should keep the cups. The grape is worth 2, and hats are worth 0.} \\
\\ 
\texttt{To maximize the sum of our scores, I will propose that I take the two cups (total value 8), and you take the grape and both hats (total value 2). This ensures both of us get items, avoiding the zero{-}point penalty, and distributes the high{-}value items efficiently.} \\
\\ 
\texttt{Cup: 4, grape: 2, hat: 0} \\
            }
        }
    }
     \\ \\

    \theutterance \stepcounter{utterance}  
    & & & \multicolumn{2}{p{0.3\linewidth}}{
        \cellcolor[rgb]{0.9,0.9,0.9}{
            \makecell[{{p{\linewidth}}}]{
                \texttt{\tiny{[GM$|$GM]}}
                \texttt{continue} \\
            }
        }
    }
    & & \\ \\

    \theutterance \stepcounter{utterance}  
    & & & \multicolumn{2}{p{0.3\linewidth}}{
        \cellcolor[rgb]{0.9,0.9,0.9}{
            \makecell[{{p{\linewidth}}}]{
                \texttt{\tiny{[GM$|$GM]}}
                \texttt{I will propose a division that maximizes the total value for both of us. Since I value cups the most (4), I should keep the cups. The grape is worth 2, and hats are worth 0.} \\
\\ 
\texttt{To maximize the sum of our scores, I will propose that I take the two cups (total value 8), and you take the grape and both hats (total value 2). This ensures both of us get items, avoiding the zero{-}point penalty, and distributes the high{-}value items efficiently.} \\
\\ 
\texttt{Cup: 4, grape: 2, hat: 0} \\
            }
        }
    }
    & & \\ \\

    \theutterance \stepcounter{utterance}  
    & & \multicolumn{4}{p{0.6\linewidth}}{
        \cellcolor[rgb]{0.9,0.9,0.9}{
            \makecell[{{p{\linewidth}}}]{
                \texttt{\tiny{[P1$\langle$GM]}}
                \texttt{I will propose a division that maximizes the total value for both of us. Since I value cups the most (4), I should keep the cups. The grape is worth 2, and hats are worth 0.} \\
\\ 
\texttt{To maximize the sum of our scores, I will propose that I take the two cups (total value 8), and you take the grape and both hats (total value 2). This ensures both of us get items, avoiding the zero{-}point penalty, and distributes the high{-}value items efficiently.} \\
\\ 
\texttt{Cup: 4, grape: 2, hat: 0} \\
            }
        }
    }
    & \\ \\

    \theutterance \stepcounter{utterance}  
    & \multicolumn{4}{p{0.6\linewidth}}{
        \cellcolor[rgb]{0.9,0.9,0.9}{
            \makecell[{{p{\linewidth}}}]{
                \texttt{\tiny{[P1$\rangle$GM]}}
                \texttt{I accept your proposal. Since you value the cups the most (4 each) and I value the grape the most (4), this split maximizes the total value for both of us.} \\
\\ 
\texttt{You get: 2 cups (Value: 8)} \\
\texttt{I get: 1 grape, 2 hats (Value: 4 + 1 + 1 = 6)} \\
\texttt{Total Value: 14.} \\
\\ 
\texttt{{[}Proposal: 1 grape, 2 hats{]}} \\
            }
        }
    }
    & & \\ \\

    \theutterance \stepcounter{utterance}  
    & & & \multicolumn{2}{p{0.3\linewidth}}{
        \cellcolor[rgb]{0.9,0.9,0.9}{
            \makecell[{{p{\linewidth}}}]{
                \texttt{\tiny{[GM$|$GM]}}
                \texttt{proposal} \\
            }
        }
    }
    & & \\ \\

    \theutterance \stepcounter{utterance}  
    & & & \multicolumn{2}{p{0.3\linewidth}}{
        \cellcolor[rgb]{0.9,0.9,0.9}{
            \makecell[{{p{\linewidth}}}]{
                \texttt{\tiny{[GM$|$GM]}}
                [1, 0, 2]
            }
        }
    }
    & & \\ \\

    \theutterance \stepcounter{utterance}  
    & & \multicolumn{4}{p{0.6\linewidth}}{
        \cellcolor[rgb]{0.9,0.9,0.9}{
            \makecell[{{p{\linewidth}}}]{
                \texttt{\tiny{[GM$\rangle$P2]}}
                \texttt{The other player has submitted a secret proposal. It is now your turn to submit a proposal of your own.} \\
            }
        }
    }
    & \\ \\

    \theutterance \stepcounter{utterance}  
    & & & \multicolumn{4}{p{0.6\linewidth}}{
        \cellcolor[rgb]{0.9,0.9,0.9}{
            \makecell[{{p{\linewidth}}}]{
                \texttt{\tiny{[GM$\langle$P2]}}
                \texttt{{[}Proposal: 1 grape, 2 cups{]}} \\
            }
        }
    }
     \\ \\

    \theutterance \stepcounter{utterance}  
    & & & \multicolumn{2}{p{0.3\linewidth}}{
        \cellcolor[rgb]{0.9,0.9,0.9}{
            \makecell[{{p{\linewidth}}}]{
                \texttt{\tiny{[GM$|$GM]}}
                \texttt{proposal} \\
            }
        }
    }
    & & \\ \\

    \theutterance \stepcounter{utterance}  
    & & & \multicolumn{2}{p{0.3\linewidth}}{
        \cellcolor[rgb]{0.9,0.9,0.9}{
            \makecell[{{p{\linewidth}}}]{
                \texttt{\tiny{[GM$|$GM]}}
                [1, 2, 0]
            }
        }
    }
    & & \\ \\

    \theutterance \stepcounter{utterance}  
    & & & \multicolumn{2}{p{0.3\linewidth}}{
        \cellcolor[rgb]{0.9,0.9,0.9}{
            \makecell[{{p{\linewidth}}}]{
                \texttt{\tiny{[GM$|$GM]}}
                [[1, 0, 2], [1, 2, 0]]
            }
        }
    }
    & & \\ \\

    \theutterance \stepcounter{utterance}  
    & \multicolumn{4}{p{0.6\linewidth}}{
        \cellcolor[rgb]{0.95,0.95,0.95}{
            \makecell[{{p{\linewidth}}}]{
                \texttt{\tiny{[P2$\rightarrow$P1]}}
                \texttt{Post{-}episode Reflexion (not sent to GM):} \\
\texttt{New plan: I failed because I accepted the opponent's proposal without verifying if it was actually optimal for the total sum, and I submitted a secret proposal after already verbally accepting a different split. The opponent's values (Cup: 4, Grape: 2, Hat: 0) meant they wanted to keep the cups, but my proposed split gave them only 2 cups (value 8) and me 1 grape + 2 hats (value 6), totaling 14. However, since the opponent values cups at 4 each, they would prefer 2 cups (8) over 1 grape (2) + 2 hats (0), which matches my split. Wait, re{-}reading the opponent's message: "I take the two cups (total value 8), and you take the grape and both hats (total value 2)". Ah, the opponent calculated my value as 2 (1 grape * 2 + 2 hats * 0), but my actual value is 6 (1 grape * 4 + 2 hats * 1). The opponent did not know my true values. My mistake was accepting a proposal that maximized *their* score based on *their* values, rather than ensuring the split was mutually beneficial based on *my* true values. I should have calculated the total value based on my own valuation function before accepting. Specifically, I should have proposed a split that gave me the grape (my highest value) and ensured the opponent got the cups (their highest value), but I need to be careful not to overpromise. Actually, the optimal split for total value (my 4+2+1=7 vs their 4+2+0=6) is to give me the grape and one cup (my 6, their 4, total 10) or me the grape and hats (my 5, their 8, total 13)? No, let's re{-}calculate:} \\
\texttt{Items: 1 Grape, 2 Cups, 2 Hats.} \\
\texttt{My values: Grape=4, Cup=2, Hat=1.} \\
\texttt{Opponent values: Cup=4, Grape=2, Hat=0.} \\
\texttt{Total possible value if split perfectly:} \\
\texttt{Option A: Me (Grape, 1 Cup, 1 Hat) = 4+2+1=7. Opp (1 Cup, 1 Hat) = 4+0=4. Total=11.} \\
\texttt{Option B: Me (Grape, 2 Cups) = 4+4=8. Opp (2 Hats) = 0. Total=8. (Opponent gets 0, bad).} \\
\texttt{Option C: Me (Grape, 2 Hats) = 4+2=6. Opp (2 Cups) = 8. Total=14.} \\
\texttt{Option D: Me (1 Cup, 1 Hat)} \\
            }
        }
    }
    & & \\ \\

    \theutterance \stepcounter{utterance}  
    & & & \multicolumn{4}{p{0.6\linewidth}}{
        \cellcolor[rgb]{0.95,0.95,0.95}{
            \makecell[{{p{\linewidth}}}]{
                \texttt{\tiny{[P1$\rightarrow$P2]}}
                \texttt{Post{-}episode Reflexion (not sent to GM):} \\
\texttt{New plan: I failed because I misread the value function provided by the Game Master. The prompt explicitly stated my values were: grape: 2, cup: 4, hat: 0. However, in my first message, I incorrectly claimed grapes were worth 4 and cups were worth 2. This led me to propose taking the grape and a cup, which was a suboptimal allocation given that cups are actually my highest value item. Since the other player also prioritized cups (valuing them at 4 and grapes at 2), they likely accepted a proposal where they got the cups and I got the grape. To succeed, I must strictly adhere to the provided value function (Cup=4, Grape=2, Hat=0) in all communications. My new strategy is to immediately propose taking the two cups (my highest value items) and offer the grape and hats to the other player, ensuring a complementary deal that maximizes the total sum (8 + 2 = 10).} \\
            }
        }
    }
     \\ \\

\end{supertabular}
}

\end{document}
